% Part: applied-modal-logics
% Chapter: epistemic-logic
% Section: properties-accessibility

\documentclass[../../../include/open-logic-section]{subfiles}

\begin{document}

\olfileid[hi]{aml}{el}{acc}

\olsection{अभिगम्यता संबंध और ज्ञानात्मक सिद्धांत}

सामान्य मोडल तर्कशास्त्र में फ़्रेम-अनुरूपता के बारे में हम जो जानते हैं,
उसके आधार पर यह देखना उपयोगी होगा कि ज्ञानात्मक व्याख्या देने पर अभिलाक्षणिक
!!{formula} कैसे दिखते हैं। हम पहले ही कह चुके हैं कि ज्ञानात्मक तर्कशास्त्रों
की व्याख्या सामान्यतः S5 में की जाती है। अतः अब देखें कि विभिन्न ज्ञानात्मक
सिद्धांत किस प्रकार निरूपित होते हैं और वे अलग-अलग फ़्रेम-शर्तों के अनुरूप
कैसे हैं।

सामान्य मोडल तर्कशास्त्र से स्मरण करें कि भिन्न मोडल !!{formula} अभिगम्यता
संबंधों के भिन्न गुणों को अभिलक्षित करते थे। निम्न सारणी उनमें से कुछ को
चुनती है जो विशिष्ट ज्ञानात्मक सिद्धांतों के अनुरूप हैं।

\begin{table}[t]
    \begin{tabular}{| p{.48\textwidth} || p{.48\textwidth} |}
      \hline
      {\emph{यदि $R$ \dots है}} & {\emph{तो \dots~$\mModel{M}$ में सत्य है:}} \\
      \hline \hline
        
      & $\Knows (p \lif q) \lif (\Knows p \lif \Knows q)$ 
      \hfill \newline{(संवृतता)} \\
      \hline
      \emph{स्वपरावर्ती}: $\forall w Rww$  
      & $\Knows p \lif p$ \hfill \newline{(सत्यपरकता)} \\
      \hline
      \emph{संक्रमणशील}: \newline
      $\forall u \forall v \forall w ((Ruv \land Rvw) \lif Ruw)$ & 
      $\Knows p \lif \Knows \Knows p$ \hfill 
           \newline{(सकारात्मक आत्मनिरीक्षण)} \\
      \hline 
      \emph{यूक्लिडीय}: \newline
      $\forall w \forall u \forall v ((Rwu \land Rwv) \lif Ruv)$ &  
      $\lnot \Knows p \lif \Knows \neg \Knows p$ \hfill 
        \newline{(नकारात्मक आत्मनिरीक्षण)} \\
      \hline
    \end{tabular}
    \caption{चार ज्ञानात्मक सिद्धांत।}
    \ollabel{tab:four}
  \end{table} 

$T$ अभिगृहीत के अनुरूप सत्यपरकता को प्रायः इन सिद्धांतों में सबसे कम
विवादास्पद माना जाता है, क्योंकि वह यह दावा व्यक्त करती है कि यदि कोई
!!{formula} ज्ञात है, तो उसे सत्य होना चाहिए। संवृतता और सकारात्मक तथा
नकारात्मक आत्मनिरीक्षण कहीं अधिक विवादास्पद हैं।

$K$ अभिगृहीत के अनुरूप संवृतता यह विचार व्यक्त करती है कि किसी कर्ता का
ज्ञान निहितार्थ के अधीन संवृत है। कुछ मामलों में यह विश्वसनीय लग सकता है।
मसलन, मुझे पता हो सकता है कि यदि मैं विक्टोरिया में हूँ तो मैं वैंकूवर द्वीप
पर हूँ। विचित्र संशयवादी परिदृश्यों को छोड़कर, मैं यह जानता हूँ कि मैं
विक्टोरिया में हूँ; इससे यह भी निकलना चाहिए कि मैं जानता हूँ कि मैं वैंकूवर
द्वीप पर हूँ। इस मामले में मेरे ज्ञान का तार्किक संवरण अपेक्षाकृत सहज लग
सकता है। दूसरी ओर, हम अपने ज्ञान के सभी परिणामों पर सदैव विचार नहीं करते;
अतः दूसरे मामलों में इससे कम सहज निष्कर्ष मिल सकते हैं।

सकारात्मक आत्मनिरीक्षण, जिसे कभी-कभी KK-सिद्धांत कहते हैं, प्रायः इस कथन
के रूप में व्यक्त किया जाता है कि यदि मैं कुछ जानता हूँ, तो मैं यह भी जानता
हूँ कि मैं उसे जानता हूँ। यह अभिगृहीत 4 का ज्ञानात्मक प्रतिरूप है। इसी के
अनुरूप नकारात्मक आत्मनिरीक्षण इस कथन के रूप में व्यक्त किया जाता है कि यदि
मैं कोई बात \emph{नहीं} जानता, तो मैं जानता हूँ कि मैं उसे नहीं जानता; यह
अभिगृहीत 5 का प्रतिरूप है। इन दोनों के अपेक्षाकृत सामान्य प्रतिदृष्टांत
दिखते हैं, जिनमें मुझे इस बात का निश्चय नहीं होता कि मैं कोई ऐसी बात जानता
हूँ या नहीं जिसे वास्तव में मैं जानता हूँ।


\end{document}
